\section{The question and main finding}
\label{sec:question}

A model receives a short memory stating where someone lives. Asked for the city, it writes a note quoting the
relevant fact, then answers ``unknown''. This is one form of \emph{false abstention}: the context supports an
answer, but the model declines to give it. An appropriate abstention is different---the context really does not
supply enough evidence. A useful intervention must recover the missed answers while preserving that distinction.
The note is a prompted summary, not evidence that the model completed a reasoning process
(\cref{sec:setup}).

This question grew out of a project on memory management. The broader story appears in
\href{https://qamil-mirza.com/research/more-than-reasoning}{\emph{More Than Reasoning}} and
\href{https://qamil-mirza.com/research/ai-brain-surgery}{\emph{Weight-Space Programming for Mechanistic Interpretability}}.
Here we focus on one experiment: can an AI agent propose a small weight edit that improves answering on examples
it has never seen, and does that edit offer a better trade-off than simpler interventions?

\paragraph{What motivated the experiment.}
Earlier experiments revealed a gap between recovering an answer-related signal and explaining a failure.
On Qwen2.5-3B-Instruct, the gold answer's first token appeared in the final-layer top-10 on 0.64 of failures;
the corresponding rates were 0.68--0.80 in three other model families. This is a readout-qualified observation,
not proof that the answer was derived from the supplied evidence. On controlled tasks, vocabulary lenses
recovered answers only in the last six layers. That timing reflects the readout's limits as well as the
model's computation.

The model's own generated notes made the problem especially visible. Of 165 held-out false abstentions,
152 named the gold city before answering ``unknown''. The city was written down; that alone does not show
that the model knew the answer. Asking the model whether the memory was sufficient did not resolve the issue:
Qwen 3B answered ``no'' on most rows it could answer correctly.

Instructions and activation steering could recover answers. A content-free direction added at layer 29
encouraged answering, following the refusal-direction line of work
\citep{arditi2024refusal,lavi2025unanswerability}. The evaluated ungated interventions also increased unsupported
answers on unanswerable questions: rates reached 0.188 for the layer-29 push and 0.451 for removing the
``unknown'' option. A separate, template-limited gating result qualified this pattern; it did not establish a
general solution. Some earlier observations replicated known results, and the steering work included
exploratory and secondary analyses.

False abstention was more frequent in Qwen 3B under the tested conditions: its in-view refusal rate was 0.30,
compared with 0.10 in Llama 3B, 0.02 in Gemma 4B and 0.03 in Phi-4-mini, which more often answered incorrectly
instead. This is a frequency contrast, not a claim that the behaviour is unique to Qwen.

\paragraph{Main finding.}
An agent-proposed edit to the final attention layer met the pre-registered repair criterion on sealed data.
Its trade-off between recovering correct answers and producing unsupported ones was better than an added prompt
instruction, an activation-steering intervention and directional orthogonalisation. Superiority over a logit bias
against ``unknown'', or over removing that option from the prompt, was not established. The following sections
describe what the agent could change, how the edits were selected, and what the independent tests support.
